\documentclass[12pt]{jlreq}
\usepackage{amsmath, amssymb}

\newcommand{\C}{\mathbb{C}}
\newcommand{\R}{\mathbb{R}}
\newcommand{\Z}{\mathbb{Z}}
\newcommand{\Tr}{\operatorname{Tr}}
\newcommand{\Impart}{\operatorname{Im}}
\newcommand{\AF}{\operatorname{AF}}
\newcommand{\AX}{\operatorname{AX}}
\newcommand{\GL}{\operatorname{GL}}

\begin{document}

区間 \((0,1)\) 内に値をとる \(C^\infty\) 級関数
\(u=u(t,x),\ t>0,\ x\in\R\) が与えられ，任意の \(t>0\) に対し積分
\[
  \int_x^\infty u(t,z)\,dz
\]
が収束するとする。また，任意の \(T>1\) に対して \(\R\) 上の可積分関数 \(\varphi_T(x)\) が存在し
\[
  \left|\frac{\partial u}{\partial t}(t,x)\right|\le \varphi_T(x),
  \qquad
  t\in\left[\frac1T,T\right],\ x\in\R
\]
を満たすものとする。さらに
\[
  f_t(x)=x+\int_x^\infty u(t,z)\,dz,\qquad x\in\R
\]
は \(\R\) から \(\R\) への全単射を定めるものとする。関数 \(v=v(t,y)\) を
\[
  v(t,y)=\frac{u(t,f_t^{-1}(y))}{1-u(t,f_t^{-1}(y))},
  \qquad t>0,\ y\in\R
\]
で定めるとき，次の問いに答えよ。ただし \(f_t^{-1}\) は \(f_t\) の逆写像を表す。

\begin{enumerate}
  \item \(u\) が熱方程式
  \[
    \frac{\partial u}{\partial t}(t,x)=\frac{\partial^2u}{\partial x^2}(t,x),
    \qquad t>0,\ x\in\R
  \]
  および
  \[
    \lim_{x\to\infty}\frac{\partial u}{\partial x}(t,x)=0
  \]
  を満たすとき，\(v\) は方程式
  \[
    \frac{\partial v}{\partial t}(t,y)
      =\frac{\partial^2}{\partial y^2}\left\{\frac{v(t,y)}{1+v(t,y)}\right\},
    \qquad t>0,\ y\in\R
  \]
  を満たすことを示せ。
  \item \(u\) が非粘性 Burgers 方程式
  \[
    \frac{\partial u}{\partial t}(t,x)
      =\frac{\partial}{\partial x}\{u(t,x)(1-u(t,x))\},
    \qquad t>0,\ x\in\R
  \]
  および
  \[
    \lim_{x\to\infty}u(t,x)(1-u(t,x))=0
  \]
  を満たすとき，\(v\) が満たす \(1\) 階の偏微分方程式を求めよ。
\end{enumerate}

\end{document}